% SEGMENT OLP-0678-S01
% Part: proof-theory
% Chapter: normalization
% Section: translations

\documentclass[../../../include/open-logic-section]{subfiles}

\begin{document}

\olfileid{pt}{nor}{tr}

\olsection{இயல்பு வடிவ \Log{N2i}-\usetoken{P}{proof} களுக்கும் \Log{G2i} க்கும் இடையிலான மாற்றம்}

\Log{N2i} இல் உள்ள !!^{proof}s ஐ \Log{G2i} இல் உள்ள
!!{proof}s ஆகவும் மறுதிசையிலும் மாற்றலாம். முதலில்,
வெட்டில்லாத \Log{G2i} !!{proof}s ஐ மாற்றும்போது
இயல்பு வடிவ \Log{N2i}-!!{proof}s கிடைக்கும்.

இந்த முடிவைச் சொல்ல, \Log{N2i}, \Log{G2i}
ஆகியவற்றின் தொடரணி இடப்புறங்களை ஒப்பிட ஒரு சொல்
தேவை. அடையாளமிட்ட வாய்ப்பாடுகளின் கணம்~$\Gamma'$,
அடையாளமற்ற !!{formula}s இன் பன்மைக் கணம்~$\Gamma$
க்கு \emph{ஒத்திருக்கிறது} என்பதன் பொருள்:
ஒவ்வொரு !!{formula}~$!A$ க்கும்,
$\Gamma'$ இல் $x : !A$ வடிவில் உள்ள !!{element}s
இன் எண்ணிக்கை, $\Gamma$ இல் $!A$ இன்
நிகழ்வெண்ணை மீறக்கூடாது.

\begin{cor}
$\Log{G2i} \Proves \Gamma \Sequent \Delta$ எனில்,
$\Gamma' \Sequent !C$ க்கு இயல்பு வடிவ
\Log{N2i}-!!{proof}~$\delta$ உண்டு. இங்கே
$\Delta$ வெற்றானால் $!C \ident \lfalse$;
$\Delta=\{!A\}$ என்றால் $!C$ அந்த ஒற்றை~$!A$.
மேலும் அடையாளமிட்ட !!{formula}s இன் கணம்~$\Gamma'$,
அடையாளமற்ற !!{formula}s இன் பன்மைக் கணம்~$\Gamma$
க்கு ஒத்திருக்கிறது; அதாவது ஒவ்வொரு
!!{formula}~$!A$ க்கும், $\Gamma'$ இல்
$x : !A$ வடிவில் உள்ள !!{element}s இன் எண்ணிக்கை,
$\Gamma$ இல் $!A$ இன் பிரதிகளின் எண்ணிக்கையை
மீறாது.
\end{cor}

\begin{proof}
\olref[nat][tgi]{prop:G2i-to-N2i} இன்
நிறுவலைப் பரிசீலித்தால் இது தெரிகிறது:
!!{introduction} விதிகளை !!{proof}s இன்
இறுதியில் சேர்க்கிறோம்; !!{elimination} விதிகளை
இத்தகைய !!{proof}s இன் மேற்பகுதியில் மட்டுமே சேர்க்கிறோம்.
ஆகவே !!a{elimination} விதிக்கு மேலாக
!!{introduction} விதியை உருவாக்குவதில்லை.
\end{proof}

% SEGMENT OLP-0678-S02
ஆனால் \olref[nat][tgi]{prop:G2i-to-N2i} இல்
தரப்பட்ட \Cut{} விதியின் மாற்றம், கிடைக்கும்
\Log{N2i}-!!{proof}~$\delta$ இல் வெட்டுகளை
உருவாக்கலாம். \Cut{} விதியுடன் விரிவாக்கப்பட்ட
\Log{G2i} !!{proof} அந்த விதியில் முடிவதாகவும்,
அதன் உடனடி துணை-!!{proof}s $\pi_1$, $\pi_2$
முறையே $\Gamma_1 \Sequent !A$,
$!A, \Gamma_2 \Sequent \Delta_2$ ஆகியவற்றின்
நிறுவல்களாகவும் இருக்கட்டும். $\pi_2$ இன்
மாற்றமான~$\delta_2$ இல் வரும்
$x:!A \Sequent !A$ அடிகோள்களில், $\pi_1$
இன் மாற்றமான~$\delta_1$ ஐ ஒட்டுகிறோம்.
$\delta_1$ இன் கடைசி விதி, முதன்மை
!!{formula}~$!A$ ஐக் கொண்ட
!!a{introduction} விதியாகவும், $\delta_2$
இன் $x:!A \Sequent !A$ அடிகோள் ஓர்
!!a{elimination} விதியின் முதன்முன்கோளாகவும்
இருந்தால், விளைவான~$\delta$ இல் வெட்டு உருவாகும்.

\Log{N2i} இன் !!{proof}s ஐ~\Log{G2i}
இன் !!{proof}s ஆகவும் மாற்றலாம்.
\olref[nat][tni]{prop:N2-to-G2} இன்
நிறுவலில் \Elim{\land}, \Elim{\lif},
\Elim{\lnot}, \Elim{\lforall} ஆகியவற்றை
மாற்றும்போது கிடைக்கும் \Log{G2i} !!{proof}
இல் \Cut{} அனுமானங்கள் தோன்றுகின்றன.
ஆனால் தொடக்க \Log{N2i}-!!{proof}
இயல்பு வடிவில் இருந்தால், அவற்றைத் தவிர்க்கலாம்.

% SEGMENT OLP-0678-S03
\Log{N2i}-!!{proof}~$\delta$ இல் வரும்
தொடரணி நிகழ்வுகளின் $S_1$, \dots,~$S_n$
என்ற பட்டியலை \emph{கிளை} என்போம்:
$S_1$ ஓர் அடிகோளாகவும், $S_i$ ஓர்
அனுமானத்தின் முதன்முன்கோளாக இருக்கும்
ஒவ்வொரு நேர்விலும் $S_{i+1}$ அந்த
அனுமானத்தின் முடிவாகவும் இருக்க வேண்டும்.
அதாவது, ஒரு கிளை அடிகோளிலிருந்து~$\delta$
இல் கீழ்நோக்கித் தொடரணிகளைப் பின்தொடர்கிறது;
ஆனால் !!{elimination} விதிகளின்
சிறுமுன்கோள்களில் நின்றுவிடுகிறது.
$S_n$ இறுதித் தொடரணியாக இருக்கும் கிளை~$\delta$
இன் \emph{முதன்மைக் கிளை}. ஒவ்வொரு
!!{proof}~$\delta$ வுக்கும் குறைந்தது ஒரு
முதன்மைக் கிளை உண்டு: ஒவ்வொரு விதிக்கும்
குறைந்தது ஒரு முதன்முன்கோள் உண்டு;
\Intro{\land} க்கு மட்டுமே இரண்டு உள்ளன.

% SEGMENT OLP-0678-S04
\begin{prop}\ollabel{prop:normal-N2i-to-G2i}
$\delta$ என்பது~$\Gamma \Sequent !A$ இன்
இயல்பு வடிவ $\Log{N2c}$ அல்லது
$\Log{N2i}$-!!{proof} எனில்,
$\Gamma' \Sequent \Delta$ இன் வெட்டில்லாத
$\Log{G2c}$-!!{proof} (முறையே
$\Log{G2i}$-!!{proof})~$\pi$ உண்டு.
இங்கே $\Gamma$ இலிருந்து அடையாளங்களை
நீக்குவதால் கிடைக்கும் !!{formula}s இன்
பன்மைக் கணம்~$\Gamma'$. $!A$ என்பது
$\lfalse$ அல்ல எனில் $\Delta=\{!A\}$;
அதுவே அபத்தம் எனில் $\Delta=\emptyset$.
\end{prop}

\begin{proof}
இம்முறை $\Gamma \Sequent !A$ இன்
!!{proof}~$\delta$ விலுள்ள அனுமானங்களின்
எண்ணிக்கையின் மீது தொகுத்தறிகிறோம்.
$\delta$ இல் அனுமானமே இல்லை எனில்,
அது $x:!A \Sequent !A$ என்ற அடிகோள்.
அதற்கொத்த $!A \Sequent !A$ அடிகோளே~$\pi$;
$!A \ident \lfalse$ எனில்
$\lfalse \Sequent \ $ என்ற அடிகோள்.

$\delta$ ஓர் அனுமான விதியில்
முடிந்தால் நேர்வுகளைப் பிரிப்போம்:
\begin{enumerate}
  \item $R$ என்பது !!a{introduction} விதியோ
  \FalseInt ஓ எனில், மாற்றம்
  \olref[nat][tni]{prop:N2-to-G2} இன்
  நிறுவலில் உள்ளபடியே நடக்கிறது;
  \Cut{} தேவையில்லை.

  \item $R$ ஓர் !!a{elimination} விதி
  எனில், முதன்மைக் கிளைகள் குறித்துப்
  பின்வருவனவற்றைக் கவனிப்போம்:
  \begin{enumerate}
    \item $\delta$ இன் முதன்மைக் கிளையில்
    !!{introduction} விதியோ \FalseInt{}
    விதியோ இருக்காது. அப்படியிருந்தால்,
    அதே கிளையில் !!a{introduction} விதியையோ
    பொய்யிலிருந்து வருவித்தலையோ ஓர்
    !!a{elimination} விதி தொடர்ந்து
    இயல்பு வடிவத்துக்கு முரணாகும்.
    \item எனவே எந்த !!{introduction} விதியும்
    இல்லாத~$\delta$ வுக்கு ஒரே
    முதன்மைக் கிளை உண்டு. \Intro{\land}
    க்கு மட்டுமே இரண்டு முதன்முன்கோள்கள்
    உண்டு; பல முதன்மைக் கிளைகள் இருக்க
    அவை \Intro{\land} இன் இரு
    முன்கோள்கள் வழியே செல்ல வேண்டும்.
    \item $\delta$ வின் முதன்மைக் கிளையில் \Elim{\lor}
    அல்லது \Elim{\lexists} இன்
    முதன்முன்கோள் வந்தால், அத்தகைய விதி
    ஒன்றே இருக்கும்; அதுவே இறுதி விதி~$R$.
    இல்லையெனில் \Elim{\lor} அல்லது
    \Elim{\lexists} இன் முடிவு மற்றோர்
    !!a{elimination} விதியின்
    முதன்முன்கோளாகி, வரிசைமாற்று
    உருமாற்றத்தால் அகற்றக்கூடிய வெட்டு
    தோன்றும்.
    \item $\delta$ வின் முதன்மைக் கிளையில்
    \Elim{\lor},
    \Elim{\lexists} தவிர்ந்த
    !!{elimination} விதிகள் எடுகோள்களை
    விடுவிப்பதில்லை. \Elim{\lor},
    \Elim{\lexists} ஆகிய விதிகளும்
    சிறுமுன்கோளுக்கு மேலேயுள்ள
    எடுகோள்களை மட்டுமே விடுவிக்கின்றன.
    ஆகவே முதன்மைக் கிளையின் ஒரு தொடரணி
    $S_i$ இன் இடச்சூழல்~$\Gamma_1$
    எனில், $S_i$ ஐ முதன்முன்கோளாகக் கொண்ட
    அனுமானத்தின் முடிவான~$S_{i+1}$
    இன் இடச்சூழல் $\Gamma_1'
    \supseteq \Gamma_1$ ஆகும்.
    \item இதனால் $\delta$ வின் முதன்மைக் கிளையின்
    மேற்பகுதி தொடரணி~$S_1$
    $x: !C \Sequent !C$ எனில்,
    $x:!C \in \Gamma$.
  \end{enumerate}
  இப்போது~$!C$ இன் வடிவத்தைப் பொறுத்து
  நேர்வுகளைப் பிரிப்போம். அந்தக்
  கிளையிலுள்ள ஒவ்வொரு $S_i$ யும் ஓர்
  !!a{elimination} விதியின்
  முதன்முன்கோள்; எனவே
  $x:!C \Sequent !C$ க்கும் இது பொருந்தும்.

% SEGMENT OLP-0678-S05
  \item $!C \ident !D \land !E$ எனில்,
  $\delta$ வின் வடிவம்:
  \[
  \Axiom$x:!D \land !E \fCenter !D \land !E$
  \RightLabel{\Elim{\land}[1]}
  \UnaryInf$x:!D \land !E \fCenter !D$
  \Deduce$x:!D \land !E, \Gamma_1 \fCenter !A$
  \DisplayProof
  \]
  இங்கு முதன்மைக் கிளையில்
  $x:!D \land !E \Sequent !D$ வருகிறது.
  அந்தக் கிளையில் \Intro{\lif},
  \Intro{\lnot} அறிமுகங்களோ
  \Elim{\lor}, \Elim{\lexists}
  சிறுமுன்கோள்களோ இல்லாததால்,
  அதன் சூழலில் உள்ள !!{formula}s இல்
  $x:!D \land !E$ எடுகோள்
  விடுவிக்கப்படாமல் இறுதிச் சூழலிலும்
  இருக்கும். அதனால்தான் இறுதித்
  தொடரணியின் சூழலில்~$x:!D \land !E$
  இடம்பெறுகிறது. மேலும்,
  \Elim{\land} இல் முடியும்
  துணை-!!{proof} க்குப் பதிலாக
  $x:!D \Sequent !D$ ஐ வைத்து,
  முதன்மைக் கிளையின் ஒவ்வொரு
  தொடரணிச் சூழலிலும்
  $x:!D \land !E$ க்குப் பதிலாக
  $x:!D$ ஐ வைக்கிறோம். இவ்வாறு
  !!{proof}~$\delta$ இலிருந்து !!{proof}~$\delta_1$
  கிடைக்கும்:
  \[
  \bottomAlignProof
  \Axiom$x:!D \land !E \fCenter !D \land !E$
  \RightLabel{\Elim{\land}[1]}
  \UnaryInf$x:!D \land !E \fCenter !D$
  \Deduce$x:!D \land !E, \Gamma_1 \fCenter !A$
  \DisplayProof
  \quad\text{ இதுவாக }\quad
  \bottomAlignProof
  \Axiom$x:!D \fCenter !D$
  \Deduce$x:!D, \Gamma_1 \fCenter !A$
  \DisplayProof
  \]
  $\delta_1$ இல் $\delta$ வை விட
  $1$ அனுமானம் குறைவு; ஆகவே
  தொகுத்தறிதல் கருதுகோள் அதற்குப்
  பொருந்தும். அசல் அடையாளமிட்ட
  $x:!D,\Gamma_1 \Sequent !A$
  தொடரணியிலிருந்து, இடப்புற
  அடையாளங்களை நீக்கிய
  \Log{G2i}-!!{proof}~$\pi_1$
  கிடைக்கிறது. \LeftR{\land} ஐப்
  பயன்படுத்தி !!{proof}~$\pi$ ஐப் பெறுகிறோம்:
  \[
  \AxiomC{}
  \RightLabel{$\pi_1$}
  \Deduce$!D, \Gamma_1' \fCenter !A$
  \RightLabel{\LeftR{\land}}
  \UnaryInf$!D \land !E, \Gamma_1' \fCenter !A$
  \DisplayProof
  \]
  $!A \ident \lfalse$ எனில்,
  இந்தக் கட்டத்திலேயே முடிவின்
  வலப்புறம் வெற்றாக இருக்க வேண்டும்;
  இதுவே முன்மொழிவின்
  $\Delta=\emptyset$ நேர்வு.
  பின்வரும் மூலக் காட்சி
  \RightR{\Weakening} ஐ மேலும்
  சேர்த்து வலப்புறத்தில் அபத்தத்தை
  வைக்கிறது; அதன் இறுதி தொடரணி
  வெற்று-வலப்புற முடிவு அல்ல:
  \[
  \AxiomC{}
  \RightLabel{$\pi_1$}
  \Deduce$!D, \Gamma_1' \fCenter $
  \RightLabel{\LeftR{\land}}
  \UnaryInf$!D \land !E, \Gamma_1' \fCenter $
  \RightLabel{\RightR{\Weakening}}
  \UnaryInf$!D \land !E, \Gamma_1' \fCenter !A$
  \DisplayProof
  \]

% SEGMENT OLP-0678-S06
  \item $!C \ident !D \lor !E$ எனில்,
  அது \Elim{\lor} இன்
  முதன்முன்கோள். இந்த விதியே
  முதன்மைக் கிளையின் இறுதி விதி~$R$.
  எனவே $\delta$ வின் வடிவம்:
  \[
  \Axiom$x:!D \lor !E \fCenter !D \lor !E$
  \AxiomC{}
  \RightLabel{$\delta_1$}
  \Deduce$y:!D, \Gamma_1 \fCenter!A$
  \AxiomC{}
  \RightLabel{$\delta_2$}
  \Deduce$z:!E, \Gamma_2 \fCenter !A$
  \RightLabel{\Elim{\lor}}
  \TrinaryInf$x:!D \lor !E, \Gamma_1, \Gamma_2 \fCenter !A$
  \DisplayProof
  \]
  துணை-!!{proof}s $\delta_1$,
  $\delta_2$ மீது தொகுத்தறிதல்
  கருதுகோளைப் பயன்படுத்தி,
  முறையே $!D,\Gamma_1'
  \Sequent !A$,
  $!E,\Gamma_2' \Sequent !A$
  ஆகியவற்றின் \Log{G2i}-!!{proof}s
  $\pi_1$, $\pi_2$ ஐப் பெறுகிறோம்;
  $!A$ ஆனது $\lfalse$ எனில் இவற்றின்
  வலப்புறங்கள் வெற்றாகும்.
  $\LeftR{\lor}$ ஐப் பயன்படுத்தி
  $\pi$ ஐப் பெறுகிறோம்:
  \[
  \AxiomC{}
  \RightLabel{$\pi_1$}
  \Deduce$!D, \Gamma_1' \fCenter !A$
  \AxiomC{}
  \RightLabel{$\pi_2$}
  \Deduce$!E, \Gamma_2' \fCenter !A$
  \RightLabel{\LeftR{\lor}}
  \BinaryInf$!D \lor !E, \Gamma_1', \Gamma_2' \fCenter !A$
  \DisplayProof
  \]
  $\Elim{\lor}$ ஆனது $y: !D$
  அல்லது $z: !E$ ஐ விடுவிக்காத
  நேர்வில், அவை தொடர்புடைய
  சிறுமுன்கோள்களின் சூழலில்
  இல்லாததால், தேவையான
  \LeftR{\Weakening} ஐச்
  சேர்க்கிறோம். இறுதி வாய்ப்பாடு
  அபத்தம் எனில், முன்மொழிவுக்கான
  முடிவின் வலப்புறம் வெற்றாகவே
  இருக்க வேண்டும். பின்வரும்
  மூலக் காட்சியில் இடப்புற
  பலவீனமாக்கல்களுடன் இறுதியில்
  \RightR{\Weakening} உம்
  சேர்ந்து வலப்புறத்தில் அபத்தம்
  வருகிறது; வெற்று-வலப்புற
  முடிவைப் பெற அந்த இறுதிப்
  படியை விட வேண்டும்:
  \[
  \AxiomC{}
  \RightLabel{$\pi_1$}
  \Deduce$\Gamma_1' \fCenter $
  \RightLabel{\LeftR{\Weakening}}
  \UnaryInf$!D, \Gamma_1' \fCenter $
  \AxiomC{}
  \RightLabel{$\pi_2$}
  \Deduce$\Gamma_2' \fCenter $
  \RightLabel{\LeftR{\Weakening}}
  \UnaryInf$!E, \Gamma_2' \fCenter $
  \RightLabel{\LeftR{\lor}}
  \BinaryInf$!D \lor !E, \Gamma_1', \Gamma_2' \fCenter $
  \RightLabel{\RightR{\Weakening}}
  \UnaryInf$!D \lor !E, \Gamma_1', \Gamma_2' \fCenter !A$
  \DisplayProof
  \]

% SEGMENT OLP-0678-S07
  \item $!C \ident !D \lif !E$
  எனில், $\delta$ வின் வடிவம்:
  \[
  \Axiom$x:!D \lif !E \fCenter !D \lif !E$
  \AxiomC{}
  \RightLabel{$\delta_1$}
  \Deduce$\Gamma_1 \fCenter!D$
  \RightLabel{\Elim{\lif}}
  \BinaryInf$x:!D \lif !E, \Gamma_1 \fCenter !E$
  \RightLabel{$\delta_2$}
  \Deduce$x:!D \lif !E, \Gamma_1, \Gamma_2 \fCenter !A$
  \DisplayProof
  \]
  முதன்மைக் கிளையில்
  $x:!D \lif !E,\Gamma_1
  \Sequent !E$ உள்ளது. அந்தக்
  கிளையில் \Intro{\lif},
  \Intro{\lnot} அறிமுகங்களோ
  \Elim{\lor}, \Elim{\lexists}
  சிறுமுன்கோள்களோ இல்லாததால்,
  $x:!D \lif !E,\Gamma_1$ இன்
  சூழலில் உள்ள !!{formula}s
  இறுதிச் சூழலிலும் எஞ்சுகின்றன;
  ஆகவே அந்தச் சூழல் இறுதித் தொடரணியின்
  இடப்புறத்திலும் வருகிறது:
  $x:!D \lif !E,\Gamma_1$.
  $\delta_2$ என்ற துணை-!!{proof}
  இன் தொடக்கத்தில் \Elim{\lif}
  இல் முடியும் பகுதியை
  $x:!E \Sequent !E$ ஆல்
  மாற்றி, அதனுடைய முதன்மைக்
  கிளையின் ஒவ்வொரு தொடரணிச்
  சூழலிலும்
  $x:!D \lif !E,\Gamma_1$
  க்குப் பதிலாக $x:!E$ ஐ
  வைக்கிறோம். இதனால்
  !!{proof}~$\delta_2'$ கிடைக்கும்:
  \[
  \bottomAlignProof
  \Axiom$x:!D \lif !E, \Gamma_1 \fCenter !E$
  \RightLabel{$\delta_2$}
  \Deduce$x:!D \lif !E, \Gamma_1, \Gamma_2 \fCenter !A$
  \DisplayProof
  \quad\text{ இதுவாக }\quad
  \bottomAlignProof
  \Axiom$x:!E \fCenter !E$
  \RightLabel{$\delta_2'$}
  \Deduce$x:!E, \Gamma_2 \fCenter !A$
  \DisplayProof
  \]
  !!{proof}s ஆன $\delta_1$, $\delta_2'$
  இரண்டுக்கும் தொகுத்தறிதல்
  கருதுகோள் பொருந்தும்.
  $\delta$ வின் அனுமான எண்ணிக்கை,
  $\delta_1$ மற்றும் $\delta_2$
  வின் அனுமான எண்ணிக்கைகளின்
  கூட்டுத்தொகையைவிட $1$ அதிகம்;
  அந்த ஒன்று முதன்மைக் கிளைத்
  தொடக்கத்திலுள்ள \Elim{\lif}.
  இதுவே~$\delta$ இன் இறுதி விதி~$R$ ஆக
  இருந்தால், $\delta_1$ இல்
  அனுமானங்கள்~$0$ ஆகவும்
  இருக்கலாம்.

  தொகுத்தறிதல் கருதுகோளால்
  $\Gamma_1' \Sequent !D$
  மற்றும் $!E,\Gamma_2'
  \Sequent !A$ தொடர்பான
  \Log{G2i}-!!{proof}s
  $\pi_1$, $\pi_2$ கிடைக்கும்;
  அவற்றின் இடப்புறங்களில்
  அடையாளங்கள் நீக்கப்பட்டவை,
  முடிவு அபத்தமான
  நேர்வுகளில் வலப்புறம்
  வெற்றானது. \LeftR{\lif}
  ஐப் பயன்படுத்தி !!{proof}~$\pi$
  கிடைக்கும்:
  \[
  \AxiomC{}
  \RightLabel{$\pi_1$}
  \Deduce$\Gamma_1' \fCenter !D$
  \AxiomC{}
  \RightLabel{$\pi_2$}
  \Deduce$!E, \Gamma_2' \fCenter !A$
  \RightLabel{\LeftR{\lif}}
  \BinaryInf$!D \lif !E, \Gamma_1', \Gamma_2' \fCenter !A$
  \DisplayProof
  \]
  உட்படுத்தலின் முன்னுறுப்பு அபத்தம் எனில்,
  உட்படுத்தல்-இடது விதியின் இடது
  முன்கோளுக்கு $!D$ வலப்புறம்
  தேவைப்படுவதால், வெற்று
  வலப்புறமுள்ள முதல் நிறுவலுக்கு
  \RightR{\Weakening}
  சேர்க்கலாம். ஆனால்
  $!A \ident \lfalse$ எனில்,
  முன்மொழிவின் முடிவில்
  வலப்புறம் வெற்றாகவே இருக்க
  வேண்டும்; இரண்டாம் நிறுவலுக்கு
  வலப்புறப் பலவீனமாக்கலைச்
  சேர்க்கக் கூடாது. பின்வரும்
  மூலக் காட்சி இரண்டையும்
  சேர்க்கும் மாற்று வடிவம்;
  அதன் இறுதி வலப்புறம்
  முன்மொழிவின்
  $\Delta=\emptyset$ நேர்வைத்
  தருவதில்லை:
  \[
  \AxiomC{}
  \RightLabel{$\pi_1$}
  \Deduce$\Gamma_1' \fCenter $
  \RightLabel{\RightR{\Weakening}}
  \UnaryInf$\Gamma_1' \fCenter !D$
  \AxiomC{}
  \RightLabel{$\pi_2$}
  \Deduce$!E, \Gamma_2' \fCenter $
  \RightLabel{\RightR{\Weakening}}
  \UnaryInf$!E, \Gamma_2' \fCenter !A$
  \RightLabel{\LeftR{\lif}}
  \BinaryInf$!D \lif !E, \Gamma_1', \Gamma_2' \fCenter !A$
  \DisplayProof
  \]

\end{enumerate}
மீதமுள்ள நேர்வுகளும் இதேபோல்
கையாளப்படுகின்றன; அவை
பயிற்சிகளாக விடப்படுகின்றன.
\end{proof}

% SEGMENT OLP-0678-S08
\begin{cor}
இயல்பு வடிவ \Log{N1i} அல்லது
\Log{N2i}-!!{proof} இல் வரும்
ஒவ்வொரு !!{formula} வும், இறுதி
வாய்ப்பாட்டின் அல்லது
விடுவிக்கப்படாத எடுகோள்களின்
துணை-!!{formula} ஆகும்; தொடரணி
வடிவில் சொன்னால், இறுதித்
தொடரணியில் வரும் ஏதோ ஒரு
வாய்ப்பாட்டின்
துணை-!!{formula} ஆகும்.
\end{cor}

\begin{proof}
\olref{prop:normal-N2i-to-G2i}
இன்படி ஒவ்வொரு இயல்பு வடிவ
\Log{N2i}-!!{proof} ஐயும்
\Log{G2i}-!!{proof} ஆக
மாற்றலாம். மாற்றத்தைப்
பரிசீலித்தால்,
\Log{N2i}-!!{proof} இல்
வரும் ஒவ்வொரு !!{formula}
வும் \Log{G2i}-!!{proof}
இலும் வருகிறது.
\Log{G2i} இல் \Cut{}
இல்லாததால் அதற்கு
துணை-!!{formula} பண்பு
உண்டு. இயல்பான வருவித்தலின்
விடுவிக்கப்படாத எடுகோள்கள்
தொடரணி வடிவில் இடப்புற
சூழலாகின்றன; ஆகவே
அவற்றையும் இறுதி வாய்ப்பாட்டையும்
சேர்த்தே துணைவாய்ப்பாட்டு
உறவைச் சொல்ல வேண்டும்.
\end{proof}

\end{document}

